\begin{proof}

\begin{enumerate}[label={$\langle1\rangle$\arabic*.},leftmargin=*]
\item Fix the pair and let $U$ be the locus of its constructed
members in the base $B$.  It is open by the actual local
miniversality of \SeriesRef{push:src:E.2.1} and \SeriesRef{push:src:E.2.2}.  We prove that it is
closed without assuming a punctured analytic disc can
be selected inside an arbitrary open subset.

\item First consider any proper smooth family over a connected
Riemann surface.  If its known-pair locus is nonempty,
Ehresmann's theorem makes all its fibres diffeomorphic
to $S^6$, and that locus is open.  At any point of its closure choose a
coordinate disc meeting that locus.  The intersection
contains a nonempty open subset of the disc, so \ref{push:src:E.2.22}
makes the entire disc known.  Thus the known-pair locus
is closed and equals the connected Riemann surface.
Only this disc propagation, not a special choice of
punctured disc, is used below.

\item Let $b\in\overline U\subset B$.  If $b\in U$ there
is nothing to prove.  Otherwise choose a small local
representative and an irreducible branch $T$ through
$b$ which meets $U$ arbitrarily near $b$.  There are
only finitely many such local branches.  Its dimension
is positive.  The local parametrization theorem gives
a finite proper projection
\[
 \rho:T\longrightarrow\Delta^d,
 \qquad \rho^{-1}(0)=\{b\},
 \tag{E.2.23.1}
\]
after shrinking the representative and translating
$b$ to the origin.  This is the finite projection
of \citet[Chapter~II, Theorem~4.19, pp.~95--96]{Demailly2012E}.
The isolated central fibre is part of choosing
the generic projection for the germ; the
Weierstrass representative is then taken small
enough to contain no other point over zero.

\item Choose $a\in U\cap T$ in this representative,
with $\rho(a)\ne0$.  Let $L$ be the complex
line through $0$ and $\rho(a)$, restricted
to the polydisc.  An irreducible component
$C$ of the reduced analytic set
$\rho^{-1}(L)$ which contains $a$ has
dimension one.  Indeed $d-1$ linear
equations give local dimension at least
one, while finiteness of $\rho$ gives
dimension at most one.  Its finite
proper image in $L$ is one-dimensional,
so is all of $L$.  Therefore $C$ contains
the unique point $b$ over zero.

\item Normalize $C$.  Its normalization
$\widetilde C$ is a connected Riemann
surface, finite over $C$; analytic
normalization exists and a normal
curve is smooth by
\cite[Chapter~II, Theorems~7.6 and~7.12]{Demailly2012E}.
It has points above both $a$ and $b$.
Pull back the given proper smooth
family to $\widetilde C$.  Its
known-pair locus contains the
inverse image of the open set
$U\cap C$ near $a$.  The preceding
Riemann-surface argument makes
every fibre on $\widetilde C$
known, in particular the fibre
above $b$.  Thus $b\in U$.

\item We have proved that $U$ is open
and closed.  Since $B$ is
connected and $U$ is nonempty,
$U=B$.  A nonreduced base
does not add fibre isomorphism
classes: restriction to its
reduction preserves both the
fibres and smoothness of the
family.  Thus the same conclusion
also applies to such a base
when it is admitted in the
definition of deformation
equivalence.

\item Applying this result to each
link in a finite deformation
chain preserves the pair.
Conversely, E.1 identifies
all admissible sphere data
with a fixed pair with the
connected period family,
using its holomorphic
comparisons and its period
lattice $2kl\Z$.
The admissible half-plane
and its punctured-disc
quotient are connected.
Thus all members with
that pair are deformation
equivalent.  This proves
both directions.

\end{enumerate}
\end{proof}
